% 负数（综述）
% license CCBYSA3
% type Wiki

本文根据 CC-BY-SA 协议转载翻译自维基百科\href{https://en.wikipedia.org/wiki/Negative_number}{相关文章}

\begin{figure}[ht]
\centering
\includegraphics[width=6cm]{./figures/18d7e2ba5a51f4aa.png}
\caption{} \label{fig_FS1_1}
\end{figure}
在数学中，负数是正实数的相反数。[1]等价地说，负数是小于零的实数。负数常用于表示损失或缺少的量。例如，所欠的债务可以看作一种负资产。如果某个量（如电子的电荷）可以具有两种相反的性质，那么人们可以（有时是任意地）将其中一种记为正，另一种记为负。负数还用于描述那些可以小于零的刻度值，例如摄氏温度和华氏温度的刻度。负数的算术运算规律保证了“相反数”的直观概念能够反映在运算中。例如：$-(-3) = 3$因为“相反数的相反数”就是原数本身。

负数通常写作前面带一个负号。例如：$-3$ 表示一个大小为 3 的负量，读作“负三”或“减三”。相反地，大于零的数称为正数；零通常（但并非总是）被认为既不是正数，也不是负数。[2]正数的性质有时会用加号来强调，例如：$+3$。一般而言，数的“正负性”被称为它的符号。


除了零之外，每一个实数要么是正数，要么是负数。非负整数通常称为自然数（即 $0, 1, 2, 3, \ldots$），而正整数、负整数以及零合在一起称为整数。（有些自然数的定义并不包括零。）

在簿记中，所欠的金额常用红色数字或括号中的数字表示，作为负数的一种替代记号。负数早在中国汉代（公元前202年 – 公元220年）成书的《九章算术》中就已使用，不过其中可能包含更古老的材料。[3]大约三世纪的数学家刘徽建立了负数加减的运算规则。[4]到七世纪，印度数学家如婆罗摩笈多已经在描述负数的运用。伊斯兰数学家进一步发展了负数的减法与乘法规则，并解决了带有负系数的问题。[5]在负数概念出现之前，像丢番图这样的数学家认为负解是“错误的”，并称要求负解的方程为荒谬的。[6]西方数学家如莱布尼茨曾坚持认为负数是无效的，但在计算中仍然使用它们。[7][8]
\subsection{引言}
\subsubsection{数轴}
负数、正数与零之间的关系通常用数轴的形式来表示：
\begin{figure}[ht]
\centering
\includegraphics[width=10cm]{./figures/47011d923559daa2.png}
\caption{} \label{fig_FS1_2}
\end{figure}
数轴上越靠右的数值越大，而越靠左的数值越小。因此，零位于中间，右边是正数，左边是负数。

需要注意的是，绝对值更大的负数被认为更小。例如，尽管正数 8 大于正数 5，写作：
$$
8 > 5~
$$
但负数 8 被认为小于负数 5：
$$
-8 < -5~
$$
\subsubsection{带符号的数}
在负数的语境下，大于零的数称为正数。因此，除了零以外的每个实数要么是正的，要么是负的，而零本身被认为没有符号。正数有时会在前面加上一个加号，例如：$+3$表示正三。

由于零既不是正数，也不是负数，因此有时会使用非负来指代大于或等于零的数，而用非正来指代小于或等于零的数。零是一个中性数。
\subsubsection{作为减法的结果}
负数可以被看作是小数减大数的结果。例如，负三就是从零减去三得到的：
$$
0 - 3 = -3~
$$
一般来说，当一个较小的数减去一个较大的数时，结果是负数，并且其绝对值等于两数之差。
例如：
$$
5 - 8 = -3~
$$
因为：$8 - 5 = 3$.
\subsection{负数的日常应用}
\subsubsection{体育}
\begin{figure}[ht]
\centering
\includegraphics[width=8cm]{./figures/b61d62ef65d1aefd.png}
\caption{相对于标准杆的负高尔夫成绩。} \label{fig_FS1_3}
\end{figure}
\begin{itemize}
\item 足球与冰球：用净胜球表示进球与失球的差值；
\item 橄榄球：用净得分；
\item 板球：用净得分率；
\item 高尔夫：用相对于标准杆（par）的成绩来表示，低于标准杆为负分。
\item 冰球的正负值统计：当某个球员在场时，球队的总进球（+）减去总失球（−）的差值就是该球员的+/− 评分。球员可能拥有一个负的 +/− 评分。
\item 棒球的净得分差：如果球队的失分多于得分，则净得分差为负数。
\item 体育俱乐部积分扣除：若违反规则，俱乐部可能被扣分，从而在赛季初始时出现负的积分，直到赢得足够的分数抵消为止。[9][10]
\item 一级方程式（F1）赛车：单圈或分段时间常以与前一圈（或分段）的差值表示（可能是纪录圈，或前车刚完成的圈速）。如果更慢，则为正差值；如果更快，则为负差值。[11]
\item 田径项目（如短跑、跨栏、三级跳远和跳远）：风速助力会被测量并记录，顺风为正值，逆风为负值。[12][13]
\end{itemize}
\subsubsection{科学}
\begin{itemize}
\item 温度：低于 $0 \,^\circ\text{C}$ 或 $0 \,^\circ\text{F}$ 的温度。[14][15]
\item 地理坐标：赤道以南的纬度、以及本初子午线以西的经度记作负数。
\item 地形高度：地球表面的地形通常以海平面为基准来表示高度，可以是负值（例如：死海、死亡谷的地表高度，或泰晤士潮汐隧道的高度）。
\item 电路：当电池反向连接时，所加电压与其额定电压相反。例如，一个 6 伏电池若反向连接，则施加电压为 $-6$ 伏。
\item 离子：可以带正电或负电。
\end{itemize}
\subsubsection{金融}
\begin{itemize}
\item 财务报表：可以包含负余额，用负号或用括号标示。例如：银行账户透支、企业亏损（负收益）。[16]
\item GDP 增长率：一国国内生产总值的年增长率可能为负，这是经济衰退的一个指标。[17]
\item 通货紧缩：在少数情况下，通货膨胀率可能为负（即通货紧缩），表示物价总体下降。[18]
\item 股票与指数：如富时100指数、道琼斯指数，其日变动可能为负。
\item 融资中的负数：负数通常与债务、赤字同义，也常被称为“赤字状态”。
\item 负利率：在某些情况下，存款人需要向银行支付费用。[19][20][21]
\end{itemize}
\subsubsection{其他}
\begin{figure}[ht]
\centering
\includegraphics[width=8cm]{./figures/deacc9120853aa49.png}
\caption{电梯中的负楼层编号。} \label{fig_FS1_4}
\end{figure}
\begin{itemize}
\item 楼层编号：在建筑物中，地面以下的楼层常用负数编号。
\item 音频播放器：在便携式媒体播放器（如 iPod）播放音频文件时，屏幕显示的剩余时间可能以负数形式出现，并且随着剩余时间趋近于零而增加，其速率与已播放时间自零开始增加的速率相同。
\item 电视竞赛节目：
\item 在《QI》节目中，参赛者常常以负分结束。
\item 在《University Challenge》节目中，如果队伍抢答错误，则会被扣分，可能出现负分。
\item 在《Jeopardy!》节目中，分数以金钱形式计，若错误回答导致损失大于现有分数，则会出现负数金额。
\item 在《The Price Is Right》节目的“Buy or Sell”游戏中，如果损失金额超过当前账户中的金额，就会出现负分。
\item 政治学：在大选之间，一个政党的支持率变化被称为摇摆，可以为负。
\item 政治人物的支持率：支持率可能下降，从而显示为负的变化值。[22]
\item 电子游戏：根据模拟类型不同，负数可能表示失去生命、受到伤害、分数惩罚或资源消耗。
\item 弹性工时制度：若员工累计工时少于合同要求，则其工时表可能显示为负工时余额。此外，员工在一年中可能休假超额，未用工时会形成负余额并结转到下一年。
\item 电子键盘乐器：移调音符时，显示屏上用正数表示升高，用负数表示降低。例如：“−1” 表示下降一个半音。
\end{itemize}
\subsection{含负数的算术运算}
符号 “−” 既可以表示二元运算符（减法运算，例：$y - z$），也可以表示一元运算符（取负运算，例：$-x$，或两次出现如$-(-x)$）。一元取负运算作用于正数时是一个特殊情况，此时结果为负数（如 $-5$）。

虽然符号 “−” 的多重含义在形式上可能产生歧义，但在算术表达式中通常不会引起实际歧义，因为运算顺序规则决定了对每个 “−” 只有唯一的解释。然而，当运算符号彼此相邻出现时，可能会让人困惑或难以理解。一种解决办法是用括号将一元负号与其操作数括起来。

例如，表达式：$7 + -5$写作：$7 + (-5)$可能更清晰（尽管两者在形式上完全等价）。而减法表达式：$7 - 5$则是不同的运算，尽管其结果相同。在小学教育中，有时会在数的前面加上上标的正负号，以明确区分正数与负数，例如：
$$
-2 + -5 = -7~
$$
\subsubsection{加法}
\begin{figure}[ht]
\centering
\includegraphics[width=6cm]{./figures/531f42d342a4203d.png}
\caption{正数与负数加法的可视化表示。较大的球体代表数值绝对值较大的数。} \label{fig_FS1_5}
\end{figure}
两个负数的加法与两个正数的加法非常相似。例如：
$$
(-3) + (-5) = -8~
$$
这里的思想是：两笔债务可以合并为一笔更大的债务。

当正数与负数混合相加时，可以把负数看作是要被减去的正数。例如：
$$
8 + (-3) = 8 - 3 = 5~
$$
和
$$
(-2) + 7 = 7 - 2 = 5~
$$
在第一个例子中，8 的“资产”与 3 的“债务”相加，结果是 5 的净资产。

如果负数的绝对值更大，结果就为负数：
$$
(-8) + 3 = 3 - 8 = -5~
$$
和
$$
2 + (-7) = 2 - 7 = -5~
$$
在这里，资产小于债务，因此净结果是债务。
\subsubsection{减法}
如前所述，两个非负数相减可能会得到一个负数：
$$
5 - 8 = -3~
$$
一般来说，减去一个正数等价于加上一个相同大小的负数。因此：
$$
5 - 8 = 5 + (-8) = -3~
$$
以及：
$$
(-3) - 5 = (-3) + (-5) = -8~
$$
另一方面，减去一个负数等价于加上一个相同大小的正数。（其思想是：失去一笔债务等同于获得一笔资产。）因此：
$$
3 - (-5) = 3 + 5 = 8~
$$
以及：
$$
(-5) - (-8) = (-5) + 8 = 3~
$$
\subsubsection{乘法}
在数的乘法中，积的大小总是等于两个数的绝对值的乘积。积的符号则由以下规则决定：
\begin{itemize}
\item 一个正数与一个负数的积是负数；
\item 两个负数的积是正数。
\end{itemize}
因此：
$$
(-2) \times 3 = -6~
$$
以及：
$$
(-2) \times (-3) = 6~
$$
第一个例子的原因很简单：把三个 $-2$ 相加等于 $-6$：
$$
(-2) \times 3 = (-2) + (-2) + (-2) = -6~
$$
第二个例子的推理更复杂一些。其思想仍然是“失去一笔债务等于获得一笔资产”。在这里，失去两笔每笔 3 的债务，等价于获得 6 的资产：
$$
(-2 \,\text{debts}) \times (-3 \,\text{each}) = +6 \,\text{credit}~
$$
规定“两负相乘得正”不仅有直观解释，也是为了保证乘法遵循分配律。
例如：
$$
(-2) \times (-3) + 2 \times (-3) = (-2 + 2) \times (-3) = 0 \times (-3) = 0~
$$
由于：$2 \times (-3) = -6$因此，为了等式成立，必须有：$(-2) \times (-3) = 6$.

